% !TeX program = lualatex
% HTML (lwarp) variant of OpticsProblems.tex. Sources of the HTML pipeline live in html/; html/build-html.ps1 copies this file
% to the project root for the build and removes it (and everything lwarp makes) afterwards.
\documentclass{ProblemBookHTML}
\addbibresource{d:/Projects/LaTeX/MyPackages/bib/MyBase.bib}

\title{Оптика. Практичний курс}
\author{В.~В.~Іванова, С.~М.~Пономаренко}
\begin{document}
\maketitle
\tableofcontents
% the preface is a page of its own (prepare_html.py turns its \chapter* into \schapterpage)
\ForceHTMLPage
\htmlinput{intro}

\newcommand{\ChaptersOne}{GeomOpt,Photometry}
\newcommand{\ChaptersTwo}{Dispersion,Interference,Difraction,Polarisation}
\part{Геометрична оптика та фотометрія}
\multiinclude{\ChaptersOne}[]
\part{Фізична оптика}
\multiinclude{\ChaptersTwo}[]

% answers to the problems (written by the chapters, as in OpticsProblems.tex)
\answers
\multiinclude{\ChaptersOne}[-Answers]
\multiinclude{\ChaptersTwo}[-Answers]

% appendices (as in OpticsProblems.tex)
\appendix
\htmlinput{Additions/Spectra}
\htmlinput{Additions/Mathematics}
\htmlinput{Additions/Photometry}
\htmlinput{Additions/PhysConstants}

% bibliography: the two lists of OpticsProblems.tex under one heading
\nocite{BornWolfEng,Chizh,Kolobrodov1,Kolobrodov2,Ivanova}
\ForceHTMLPage
\schapter{Література}
\defbibheading{subheading}[\bibname]{\section*{#1}}
\addtocategory{Textbooks}{BornWolfEng,Chizh,Kolobrodov1,Kolobrodov2}
\addtocategory{Problems}{Ivanova}
\printbibliography[category=Textbooks,heading=subheading,title={Підручники та посібники}]
\printbibliography[category=Problems,heading=subheading,title={Задачники}]
\end{document}
